\section{Discussion}
\label{sec:discussion}

The expert validation result --- that 3 of 4 computationally-identified sparse clusters
are recognized by experienced designers as genuinely unexplored and commercially viable ---
provides the first empirical evidence that TIGER-guided design space exploration produces
actionable design opportunities rather than merely academic observations. The single
non-validated gap (G3: Deep-Architecture Highly-Forgiving) is informative in its own
right: three designers independently cited Terraforming Mars, revealing that the corpus
under-samples cooperative and semi-cooperative heavy games.

\subsection{Why Some Sparse Regions Are Not Opportunities}

Not all sparse regions in TIGER space correspond to viable design opportunities. Some
regions are sparse because the combination of dimension scores implies an incoherent
game experience (e.g., maximum Tension combined with maximum Range is difficult to
achieve simultaneously because accessibility requires reduced decision pressure). The
expert validation study surfaces this distinction: designers who rate a sparse cluster
as non-viable can explain \emph{why} that region is difficult to inhabit, providing
qualitative evidence that complements the quantitative clustering.

\subsection{TIGER as a Generative Tool}

The primary contribution of this paper is demonstrating that TIGER is not only
analytical (characterizing games that exist) but generative (identifying where games
could be). This positions TIGER as a complement to designer intuition rather than a
replacement: the computational analysis surfaces candidates that designers then evaluate
with their domain expertise, combining algorithmic breadth with expert judgment.

\subsection{Limitations}

The clustering approach assumes that the TIGER space is reasonably well-sampled by 150
games, but the corpus is biased toward Euro-tradition games. Sparse regions in the
corpus may be sparse because the corpus under-samples certain genres rather than
because no games exist there. The expert panel (five designers) is too small to
generalize about commercial viability with high confidence.
